% Migrated from days/day05.tex.  Paste your original problem, solution,
% remark and takeaway bodies into the slots below -- the wrappers are the only
% thing that changed.  Replace "the Day 4 result" style cross-references with
% \qref{Q004}, and recurring lemmas with \usetool{T1}.
\begin{entry}{Q005}{Two POVMs that never disagree}

\entrymeta{
  topic   = {Measurements; POVMs; operator positivity},
  est     = $\sim$10 minutes,
  origin  = Original,
  source  = {Original problem},
  prereq  = {A POVM is a family $\{E_a\}$ with $E_a \geq 0$ and $\sum_a E_a = \id$},
  status  = written-up,
}

\begin{problem}
Let $\{E_a\}_{a\in\mathcal A}$ and $\{F_b\}_{b\in\mathcal B}$ be POVMs on $\CC^d$.
Suppose there exist a state $\rho\in\Dop{\CC^d}$ and outcomes $a_0,b_0$ with
\[
  \Tr(E_{a_0}\rho)=1
  \qquad\text{and}\qquad
  \Tr(F_{b_0}\rho)=1 .
\]
Prove that
\[
  \Tr(E_{a_0}F_{b_0}\rho)=1 .
\]
In words: if two measurements each yield a deterministic outcome on $\rho$, their joint
statistics are deterministic as well --- even though $E_{a_0}$ and $F_{b_0}$ need not
commute.
\end{problem}

% ============================================================== solutions ===

\begin{solution}
Adding the identity $\Tr(\sum_a E_a \rho)=\Tr(\sum_b F_b \rho)=1$, together with the positivity of each $E_a,F_b$, leads to $$\Tr(E_a\rho)=\Tr(F_b\rho)=0,\ \forall a,b\neq a_0,b_0.$$ That means that all of them can only act on the $\ker(\rho)$ (its null subspace), which is, after diagonalizing $\rho= \sum_i \rho_i \proj{\phi_i}$ with $\rho_i>0$ (not taking into account the kernel basis $\rho_i=0$), we have that $\Tr(E_a\proj{\phi_i})=\Tr(F_b\proj{\phi_i})=0,\ \forall a,b\neq a_0,b_0$ for each $i$. Moreover,
$$E_a\proj{\phi_i}=F_b\proj{\phi_i}=0,\ \forall a,b\neq a_0,b_0$$
due to POVM consisting of positives operator and $\Tr(X\proj{\phi})=\bra{\phi}X\ket{\phi}$. Thus
$$E_a\rho=F_b\rho=0,\ \forall a,b\neq a_0,b_0$$ as well as 
$$\rho E_a=\rho F_b=0,\ \forall a,b\neq a_0,b_0$$
since they are hermitian operators (just apply the Hermitian conjugate)

Therefore
\begin{equation}
\begin{split}
     1&= \Tr(F_{b_0}\rho)\\
     &= \Tr(F_{b_0}\rho\sum_aE_a)\\
     &= \Tr(F_{b_0}\rho E_{a_0})\\
     &= \Tr(E_{a_0}F_{b_0}\rho )
\end{split}
\end{equation}
\end{solution}

\begin{takeaway}

\end{takeaway}

\attempt{
  date       = 2026-09-11,
  time       = 20 min,
  difficulty = 1,
  howitwent  = ,
}

\end{entry}
